% ------------------------------------------------------------------------
%
\begin{frame}{Abstract Syntax Trees (ASTs)}

  \begin{itemize}

    \item Depending on the context, we may use the arborescent
      depiction of terms to bring to the fore certain aspects of a
      computation.

    \item For example, it may be interesting to show how parts of the
      output (the right\hyp{}hand side) are actually \textbf{shared}
      with the input (the left\hyp{}hand side).

    \item In other words, how much (new) data is created by a given
      rule.

    \item This concept supposes that terms reside in some sort of
      space and that they can be referred to from different other
      terms.

    \item This abstract space serves as a model of an interpreter's
      memory, called the \textbf{heap}.

  \end{itemize}

\end{frame}

% ------------------------------------------------------------------------
%
\begin{frame}{Data sharing}

  Consider for instance
  \begin{center}
    \includegraphics[bb=71 656 302 721]{cat_dag}
  \end{center}
  which is the same definition of \fun{cat} as given above:
  \begin{equation*}
    \fun{cat}(\el,t) \xrightarrow{\alpha} t
    \qquad\qquad\qquad
    \fun{cat}(\cons{x}{s},t) \xrightarrow{\beta}
    \cons{x}{\fun{cat}(s,t)}
  \end{equation*}
  \begin{itemize}

    \item The arrows on certain edges denote data sharing.

    \item When trees are used to visualise terms, they are called
      \textbf{abstract syntax trees} (AST).

    \item When some trees share subtrees, the whole \textbf{forest}
      (set of trees) is called a \textbf{directed acyclic graph}
      (DAG).

  \end{itemize}

\end{frame}

% ------------------------------------------------------------------------
%
\begin{frame}{Data sharing and evaluation}

  The evaluation of \(\fun{cat}([1;2],[3])\) is:
  \begin{center}
    \includegraphics[bb=71 621 276 733]{cat123_push}
  \end{center}

  Note that each right-hand side shows only the rewritten call, not
  the whole term.

\end{frame}

% ------------------------------------------------------------------------
%
\begin{frame}{Data sharing and evaluation}

  Therefore, to reconstruct the value, we need to work leftwards and
  replace the rewritten call by the right-hand side:
  \begin{center}
    \includegraphics[bb=71 623 248 717]{cat123_pop0}
  \end{center}

  We do not show the calls that have been replaced by their values.

\end{frame}

% ------------------------------------------------------------------------
%
\begin{frame}{Data sharing and evaluation}

  \begin{center}
    \includegraphics[bb=71 621 212 715]{cat123_pop1}
    \qquad\quad
    \includegraphics[bb=71 621 212 715]{cat123_pop2}
  \end{center}

\end{frame}

% ------------------------------------------------------------------------
%
\begin{frame}{Call stack}

  We can see that the value of the second argument, [3], is shared
  with the value of the call, as well as the integers~1 and~2.

  \bigskip

  We saw how rewrites accumulated rightwards until the result had to
  be reconstructed by replacing leftwards the right-hand sides.

  \bigskip

  This dynamic is that of a stack, but a stack holding suspended
  function calls (that is, awaiting their values) and references to
  ASTs in the heap. This is the \textbf{call stack}.

  \bigskip

  The top of the stack (here, the rightmost abstract syntax tree) is
  the next term to be rewritten, unless the call on the left\hyp{}hand
  side is the original call, in which case the top is its value.

\end{frame}

% ------------------------------------------------------------------------
%
\begin{frame}{Garbage collection}

  \begin{itemize}

    \item In the abstract syntax trees (AST), the nodes with the calls
      that are replaced by their values have been immediately
      removed. In fact, they could have been removed later, in any
      order.

    \item The process which identifies and discards useless ASTs is
      called the \textbf{garbage collector} (GC).

    \item A term is useless if it cannot be accessed from anywhere in
      the call stack.

    \item Contrary to the heap, which can be thought of as an
      unordered set of trees (a forest of DAGs), the call stack is an
      ordered set that can be efficiently extended or reduced.

    \item The garbage collector can be conceived as a process
      interleaved with the process of evaluation.

    \item Like evaluators of functional programs, garbage collectors
      may also have different strategies (as to when dispose of
      useless data).
  \end{itemize}

\end{frame}
